\section{近期研究（2026 年 8--11 月）}
\paragraph{DesignAgent3D：用交互规划与几何定位约束室内物体编辑}
2026.08~\cite{2608_21438}。DesignAgent3D 先以对话把含糊要求整理为明确编辑方案，再借 SfM 几何引导分割选定三维目标。编辑阶段用纹理参考与几何约束重绘选定视图，然后以替换视图和重训练的方式把改变写回场景表示，使编辑可从新视角观察。

论文同时测试 NeRF 与 3DGS 底座。对比编辑质量时，基线也收到已澄清的提示词，3DGS 方案在所测任务的 CLIP 方向相似度均值为 0.093。这个结果对应论文测试的室内场景与语义评分，比较本身不衡量基线澄清用户意图的能力。

\paragraph{CoGeo-GS：联合语义选取、深度补全与外观细化}
2026.08~\cite{2608_26656}。CoGeo-GS 将文本驱动的分割掩码蒸馏到高斯中，为多个目标分配语义身份。移除对象后，以对齐尺度的单目深度约束扩散补全，再用几何正则化外观细化缓和跨视图与边界不一致。

作者在 Mip-NeRF 360 和 SPIn-NeRF 上分别验证单对象和多对象移除，使用给定相机内参与 COLMAP 点云初始化，并与 InFusion、SPIn-NeRF、GaussianEditor 和 Gaussian Grouping 比较。被遮挡区域只能由条件补全，论文的视觉与重建指标针对该数据和相机设置。

\paragraph{LightBridge：将二维重光照提示传播到完整高斯场景}
2026.09~\cite{2609_02543}。LightBridge 先用学习到的光照变化潜空间在一步内预测重光照方向，再将二维提示传到可见高斯并沿三维表示传播到输入轨迹未覆盖的部分。模型固定高斯几何，只预测新的球谐颜色参数，避免为每个场景反复烘焙重光照视频。

论文分别评价原相机轨迹上的视频和最终高斯资产的新视角渲染。对只输出视频的基线，作者固定输入高斯几何并微调球谐系数 5000 步后比较。完全不可见的光源无法从输入掩码直接定位，静态合成室内训练的领域覆盖也限制了对新材质和灯具的泛化。

\paragraph{高斯调色板重参数化：在视图空间编辑颜色和亮度}
2026.09~\cite{2609_03897}。这项工作将预训练高斯的颜色重参数化为球谐调色板权重，渲染时混合出逐像素权重。高阶系数采用秩一分解，并以视图空间约束鼓励权重稀疏和颜色一致。交互时在固定高斯与权重下求解调色板与每色调曲线，用同一优化接收像素、调色板或曲线约束。

论文分别评估新视角重建保真度、编辑局部性和跨视图一致性，并与直接操作原语的像素编辑比较。结果可烘焙回标准 3DGS。由于共享调色板是控制界面，相同颜色的不同物体仍可能需要额外选择信息才能独立编辑。

\paragraph{三维编辑分解基准：用四个维度区分编辑成功与副作用}
2026.09~\cite{2609_14899}。该基准将指令遵循、目标局部性、跨视图一致性和非目标内容保持分别计分。八个代表性编辑器共享相机标定、原图、指令与目标标注，同类表示还共享初始重建，结果在受控留出视图上汇总。

作者在 240 组场景和编辑任务中发现，不同指标的领先方法不同。InterGSEdit 的指令保真度领先，而直接编辑代理的方案在三项保持导向指标领先。论文明确指出指标受视觉语言、分割和深度对应误差影响，静态、有空间目标的测试也不足以评估全局重光照或动态编辑。

\paragraph{φ-RIE：从静态高斯场景构造可操控的仿真资产}
2026.09~\cite{2609_26795}。φ-RIE 利用同一对象证据关联生成候选、资产对齐、原对象高斯移除与露出背景补全，将选定对象变成仿真中能移动的资产并保留场景其他部分。物理引擎输出的对象位姿继续驱动高斯外观渲染。

论文在 50 个 ScanNet++ 场景测试 1871 个对象请求。在固定保留与匹配数量的构造比较中，证据筛选与注册重试将作者报告的 F1 从 0.336 提高至 0.383，同时增加场景构造时间。严格验收后仍有匹配资产出现几何塌陷，说明通过筛选不保证几何正确或物理有效。

\paragraph{WINGS：利用三维生成先验修补删除对象后的高斯空缺}
2026.09~\cite{2609_37816}。WINGS 在用户给定的遮罩区域移除现有高斯，并以保留区域作为条件，借助预训练 TRELLIS 的结构潜变量模型生成空缺区域的原语。它处理的是无参考的条件生成，因此允许多种与现有场景相容的外观和几何。

论文在 COR-NeRF、Inpaint360 和 360-USID 上用 FID、KID、目标裁剪区域 FID、LPIPS 及用户偏好评价。一些复杂纹理在潜空间里偏向轴对齐，朝向不同的图案较难补全；论文将此列为结构旋转处理仍需改进的情形。

\paragraph{以可信关键帧锚定文本驱动的三维高斯编辑}
2026.10~\cite{2610_08179}。文本编辑若独立处理每个视图，较差视图的纹理可能将高斯场拉向不一致的结果。View Matters 先估计各视图的编辑重要性，偏向几何可见、语义明确且与指令相关的关键帧。编辑信息从关键帧传播到相邻视图，同时避免低可信分支的噪声反向传播，随后在高斯优化中延续该权重。

论文以 23 组场景和提示词的 CLIP 对齐及方向相似度比较编辑，并通过相邻视图与消融检查一致性。这类分数不能替代细部几何或边界真值。方法仍沿用冻结的二维编辑器，因此很大的几何改动和细粒度局部结构可能受到编辑器能力限制。

\paragraph{3DTexMOR：在共用纹理空间修补移除物体后的区域}
2026.10~\cite{2610_11198}。一次移除多个物体后，各视角暴露的遮挡区域并不一致。3DTexMOR 把跨视角的内在外观观测先聚合到同一个纹理空间，在该空间补全，再把结果投回各相机为高斯补全提供共用监督。附加的几何约束引导补全区域获得合理三维结构。

作者在合成的漫反射与反光场景中报告定量结果，在真实场景中查看视觉表现。方法采用共用平面近似，因此曲面、形变物体和形状不规则区域可能在反投影时产生错位，真实场景也没有可供量化对照的完整遮挡面真值。

\paragraph{RxGS：在可重光照高斯表示中单独编辑漫反射反照率}
2026.10~\cite{2610_25681}。RxGS 从带位姿多视图建成 BRDF 感知的高斯表示，只优化漫反射反照率以匹配参考图风格。粗糙度、金属度、法向和环境光保持原值，因而编辑后的场景仍能在新光照下重渲染。

作者用 NeRF-Synthetic 和 NeRF-LLFF 对比跨视图几何一致性，并加入 TensorIR 与 Mip-NeRF 360 检验重光照。几何一致性与重光照是分开的实验；原文指出参与前者的风格化基线不支持后者。该设计的适用范围受输入逆渲染分解和只编辑漫反射的约束。

\paragraph{VidEdit3D：沿相机轨迹传播锚点编辑}
2026.10~\cite{2610_35756}。VidEdit3D 将校准多视图排列为平滑的短视频轨迹，从用户提供的已编辑锚点出发，以预训练图生视频模型传播语义改变。最终区域选择约束把编辑回写到目标高斯，帮助减少对背景的改动。

论文表 1 在作者测试集上报告 CLIP-D 0.1275 和 MEt3R 0.1126，并在附录分析了遮挡、细结构、掩码不准和锚点内容不完整的失败例。未选高斯的参数可保持不变，但当目标形状改变遮挡关系时，它们的渲染可见性仍可能改变。

\paragraph{ReTouchGS：把场景级颜色操作直接写入高斯资产}
2026.11~\cite{2611_23054}。ReTouchGS 自动从预训练高斯及相机条件预测颜色操作，也支持将同一场景的原渲染与单张编辑参考配对后拟合操作。解析变换先修改原语颜色，再由混合感知的轻量细化网络修正结果，最终输出几何保持不变的标准高斯场景。

评测使用 37 个场景并以每场景 16 个离轨迹视图比较场景一致的颜色目标。度量反映对给定目标的保真度。免渲染描述的是共同的高斯参数更新阶段，参考迁移的前置颜色拟合仍需要原始渲染和编辑后的图像。

\paragraph{GS-VTON-MVT：以多视图联合编辑与几何条件开展高斯试衣}
2026.11~\cite{2611_25979}。GS-VTON-MVT 在给定已校准的多视图人物图像后重建个人的高斯场景。多视图 Transformer 先联合编辑一组视图来个性化扩散模型，随后用衣物图以及当前高斯场景渲染的深度与法向控制迭代编辑，并随相机姿态调节参考信息的影响。

论文用含 60 位受试者的 3D-VTONBench，每人选择 4 个训练视图进行多视图编辑，将其余视图留作新视角评价，并比较 FID、KID、基于 LPIPS 的一致性和用户偏好。这个实验协议依赖多视图输入与已建场景，不能据此推广至单张图像试衣。
